\FloatBarrier
\section{Beamline port and validation: October 12 progress}

\subsection{Current port and validation status}

The current GEMC3 port includes all fifteen native beamline geometry and material variations: ten
standard run-mapped configurations and five detector-local configurations. It follows the active GEMC2
beamline dispatch; historical builders and disabled HD-ice daughters are not activated. All published
volumes use passive placement, and this passive system requires no sensitive digitization plugin or
optical surface.

Native geometry is published as beamline. The beamline\_cad companion contains the ten standard CAD
configurations plus DDVCS, while beamline\_transverse\_cad provides noWedges, forwardWedges, and
midpointWedges as separately selectable alternatives. Both native and CAD systems use the SQLite factory;
mesh coordinates remain in millimetres at scale one. The runtime inventory contains 51 STL files: 30
beamline meshes and seven meshes for each transverse alternative. The two apex meshes are optimized as
described below; the remaining 49 retain the original tessellation.

The semantic comparator passes all fifteen native geometry variations, fifteen material variations, and
fourteen CAD placement variations against the GEMC2 TEXT/GXML references. The ten standard runs also
match the named native and CAD entries in the GEMC2 SQLite reference. Seven identical duplicate CAD rows
at run 20000 are represented once per volume name; the DDVCS configuration is also available at its
detector-local run. This preserves the effective geometry without duplicating physical placements.

Saved construction smoke logs show completion for all fifteen native configurations and all three
transverse CAD alternatives with overlap checking enabled. These zero-event checks validate geometry
construction and mesh resolution, not particle transport or detector response. A separate default
beamline check after the apex optimization also exits normally with check\_overlaps=2, but records eight
GeomSolids1001 and nineteen GeomVol1002 exception blocks. These counts describe reported solid defects
and overlap/containment warnings, not unique defective meshes or overlapping pairs. The unchanged CAD
meshes still include surface-hole and facet-orientation diagnostics.

The current status is therefore a port with passing semantic comparisons and completed construction
checks, with residual geometry diagnostics retained for follow-up. Full-beamline transport performance,
overlap freedom, material-budget equivalence, magnetic-field operation, and full CLAS12 background
response equivalence have not been established by these checks. The isolated apex measurements below are
reported separately from this broader port validation.

\subsection{Apex shielding optimization and isolated benchmark}

The next optimization step covers two beamline apex shielding meshes. The lead fill is reduced from 2,880 to
1,440 triangles; the hub is renamed from Apex\_Shield\_PCAL\_Hub to Apex\_Shield\_Hub and reduced from
36,354 to 20,000 triangles. The two-file total falls from 39,234 to 21,440 triangles, a 45.4\% reduction.
Stored bytes fall from 1,961,868 to 1,072,168. All other beamline meshes remain unchanged. Materials,
coordinates, and passive placements are preserved in both repositories; all eleven CAD variations and eleven
run-dependent GEMC2 hub rows use the new name.

Topology-preserving quadric edge collapse keeps the lead opening and both hub components. Faces adjoining
hub edges with dihedral angles above five degrees are protected; no holes are filled and no component is
removed. Both originals and optimized meshes are closed and manifold, with no detected self-intersections or
exact zero-area triangles. Bidirectional checks include every vertex and approximately 200,000 edge/surface
samples in each direction. The largest sampled deviations are 0.0425 mm for the lead fill and 0.0209 mm for
the hub; these are sampled measurements, not proven global Hausdorff bounds. The lead volume changes by
-0.002287\%; the hub volume changes by -0.00000100\%. The original engineering meshes remain archived, and
the builder reproduces the optimized STL files byte for byte.

Table~\ref{tab:apex-meshes} summarizes the unchanged topology and reduced mesh inventory.

\begin{table}[htbp]
\centering
\caption{Apex engineering and optimized triangle counts; topology is preserved.}
\label{tab:apex-meshes}
\begin{tabular}{lrrrr}
\toprule
Mesh & Original & Optimized & Components & Genus \\
\midrule
Lead fill & 2,880 & 1,440 & 1 & 1 \\
Hub & 36,354 & 20,000 & 2 & 2 \\
\bottomrule
\end{tabular}
\end{table}

This follow-on benchmark uses one GEMC3 development executable dated October 11, 2026 and Geant4 11.4.3 on
macOS arm64. Its timings are independent of the GEMC2 LTCC measurements above and must not be compared
across engines. Only the apex pair is loaded, in air, with its original lead and carbon-steel materials.
Native beamline geometry and other CAD solids are omitted to isolate the substitution. Both cases use
FTFP\_BERT without optical processes, zero magnetic field, one thread, identical default production cuts,
and no GUI or event output. Fixed 5 GeV primaries originate at the target, uniform in solid angle from 0.25
to 3.5 degrees and over full azimuth. This aperture illuminates the forward shielding rather than reusing
the larger LTCC angular acceptance.

There are five alternating-order repetitions per case, with seeds 12345--12349 and matching LUND files for
each seed. Each electron, proton, positive-pion, and neutron run contains 200 events (1,000 events per
particle per geometry). This first beamline sample is smaller than the LTCC sample: showers in the thick
lead dominate its per-event cost. Tables report medians of the five runs. Initialization uses separate
zero-event runs with overlap checking disabled. Transport uses GEMC3's elapsed wall time since beamOn;
unlike the LTCC events-only CPU column, this is a wall-clock measure. Whole-process user plus system CPU
time and peak resident memory are also retained in the dataset. The original/optimized transport ratios are
ratios of medians, with no initialization subtraction.

Table~\ref{tab:apex-timing} gives loading, peak memory, and particle-transport measurements.

\begin{table}[htbp]
\centering
\caption{Median pair-only measurements over five repetitions. Transport entries are wall seconds
per 200 events; the final column is a ratio of medians.}
\label{tab:apex-timing}
\begin{tabular}{lrrr}
\toprule
Quantity & Original & Optimized & Original/optimized \\
\midrule
Initialization wall seconds & 1.151 & 0.970 & 1.19 \\
Loading peak RSS (MiB) & 186.09 & 153.64 & 1.21 \\
e- & 36.801 & 18.311 & 2.01 \\
proton & 9.329 & 5.174 & 1.80 \\
pi+ & 12.877 & 6.747 & 1.91 \\
neutron & 8.206 & 4.643 & 1.77 \\
\bottomrule
\end{tabular}
\end{table}

Median zero-event simulation wall time decreases by 15.7\% and loading peak resident memory decreases by
17.4\%. All 52 measured processes exit normally and print completion footers; all forty transport runs
report their 200 requested events. No Geant4 exception blocks are recorded in the pair-only loading,
transport, or separate check\_overlaps=2 construction samples. The semantic comparator passes all fifteen
native geometry/material variations and all fourteen CAD variations after the rename. Full beamline
construction checks are retained separately; remaining unoptimized meshes can still produce warnings. These
measurements establish the performance of this isolated sample, not full CLAS12 speed, overlap freedom, or
background-response equivalence. Stochastic shower histories may differ despite matched primaries and seeds.

\FloatBarrier

Figure~\ref{fig:apex-stl} compares the stored triangulations with matched cameras and bounds.

\begin{figure}[htbp]
\centering
\includegraphics[width=\textwidth]{figures/apex_stl_comparison.pdf}
\caption{Original (left) and optimized (right) apex lead-fill and hub STL surfaces. The views use
identical cameras and millimetre coordinates within each pair; edges show the stored triangulations.}
\label{fig:apex-stl}
\end{figure}
